\section*{Chapter \ref{ch:b1:p-block-13-15} --- The p-Block I: Boron, Carbon and Nitrogen Groups}

\begin{solution}{exo:b1:p-block-13-15:1}
\ce{NH4NO3}: $-$III in \ce{NH4+}, $+$V in \ce{NO3-}. \ce{N2O5}: $+$V.
\ce{HNO2}: $+$III. \ce{N2H4}: $-$II. \ce{Mg3N2}: $-$III.
\end{solution}

\begin{solution}{exo:b1:p-block-13-15:2}
\ce{BF3}: boron with three bonds and six electrons, trigonal planar.
\ce{NH3}: three bonds and a \omterm{def:b1:lewis-resonance-vsepr:lewis}{lone pair}. The \omterm{def:b1:lewis-resonance-vsepr:lewis}{lone pair} of N forms a \omterm{def:b1:lewis-resonance-vsepr:lewis-acid}{dative
bond} to B: \ce{BF3} is the \omterm{def:b1:lewis-resonance-vsepr:lewis-acid}{Lewis acid}, \ce{NH3} the \omterm{def:b1:lewis-resonance-vsepr:lewis-acid}{Lewis base}; in the
adduct both atoms are tetrahedral.
\end{solution}

\begin{solution}{exo:b1:p-block-13-15:3}
\ce{CO2 + 2NaOH -> Na2CO3 + H2O}; \ce{SiO2 + 2NaOH -> Na2SiO3 + H2O};
\ce{P4O10 + 12NaOH -> 4Na3PO4 + 6H2O}; \ce{Al2O3 + 6HCl -> 2AlCl3 + 3H2O};
\ce{Al2O3 + 2NaOH + 3H2O -> 2NaAl(OH)4}.
\end{solution}

\begin{solution}{exo:b1:p-block-13-15:4}
$\Delta_r G^\circ = 2 \times (-16.45) = \qty{-32.9}{kJ/mol}$, $\log K =
32.9/5.708 = 5.76$, $K = 10^{5.76}$.
\end{solution}

\begin{solution}{exo:b1:p-block-13-15:5}
Ring of six \ce{SiO3^2-} units: \ce{Si6O18^12-}. Double chain, per four
silicons: two with two shared corners (3 O, charge $-2$ each), two with
three ($2.5$ O, charge $-1$ each): \ce{Si4O11^6-}.
\end{solution}

\begin{solution}{exo:b1:p-block-13-15:6}
\qty{550}{kg} of \ce{Al(OH)3} is $550/78.0 = \qty{7.05}{kmol}$, giving
\qty{3.53}{kmol} of \ce{Al2O3}, \qty{360}{kg}; at \qty{92}{\percent}:
\qty{331}{kg}.
\end{solution}

\begin{solution}{exo:b1:p-block-13-15:7}
Boron has an empty orbital and accepts the \omterm{def:b1:lewis-resonance-vsepr:lewis}{lone pair} of a hydroxide ion
taken from water: \ce{B(OH)3 + 2H2O <=> B(OH)4- + H3O+}. The proton
released comes from a water molecule, not from boric acid.
\end{solution}

\begin{solution}{exo:b1:p-block-13-15:8}
N goes from $-$III to $+$II, five electrons per \ce{NH3}; each \ce{O2} takes
four. $4 \times 5 = 5 \times 4$: \ce{4NH3 + 5O2 -> 4NO + 6H2O}.
\end{solution}

\begin{solution}{exo:b1:p-block-13-15:9}
\ce{NH3 + HNO3 -> NH4NO3}. One tonne is $10^6/80.0 = \qty{1.25e4}{mol}$;
\qty{1.25e4}{mol} of ammonia to neutralise and \qty{1.25e4}{mol} to make
the nitric acid: $2.50 \times 10^4 \times 17.0 = \qty{425}{kg}$.
\end{solution}

\begin{solution}{exo:b1:p-block-13-15:10}
\omterm{def:b1:periodicity:table}{Second-period} atoms are small: their \textit{p} orbitals overlap well side
by side and form strong $\pi$ bonds, so \ce{N#N} and \ce{O=C=O} are small
stable molecules. \omterm{def:b1:periodicity:table}{Third-period} atoms are larger; their $\pi$ overlap is
poor, and they prefer more single bonds: \ce{P4} tetrahedra, the
\ce{SiO4} network.
\end{solution}

\begin{solution}{exo:b1:p-block-13-15:11}
$E^\circ(\ce{PbO2}/\ce{Pb^2+}) = 1.46 > E^\circ(\ce{Cl2}/\ce{Cl-}) = 1.36$~V:
\ce{PbO2 + 4H+ + 2Cl- -> Pb^2+ + Cl2 + 2H2O}. Lead(IV) is unstable with
respect to lead(II) (inert pair); silicon(II) is not a stable state and
\ce{SiO2} is not an \omterm{def:b1:nernst:couple}{oxidant}.
\end{solution}

\begin{solution}{exo:b1:p-block-13-15:12}
$160 \times 17.0/14.0 = \qty{194}{Mt}$ of ammonia, containing $194 \times
3.0/17.0 = \qty{34}{Mt}$ of hydrogen, made mostly from natural gas
(methane and steam).
\end{solution}

\begin{solution}{pb:b1:p-block-13-15:1}
\textbf{1.} :N\,$\equiv$\,N: --- a triple bond, very strong, with no polarity.
\textbf{2.} 0, $-$III, $+$II, $+$IV, $+$V, and $-$III/$+$V in \ce{NH4NO3}.
\textbf{3.} Bacteria on the roots of legumes, lightning, the Haber--Bosch
process.
\textbf{4.} Breaking the triple bond needs a \omterm{def:b1:elementary-steps:catalyst}{catalyst} that plants do not
have; bacteria with the enzyme nitrogenase do.
\textbf{5.} Nitrogen is reduced (0 to $-$III), hydrogen oxidised (0 to $+$I).
\textbf{6.} \ce{N2 + 3H2 <=> 2NH3}.
\textbf{7.} $K = 10^{2 \times 16.45/5.708} = 10^{5.76}$.
\textbf{8.} At \qty{25}{\celsius} the reaction is far too slow; the
\omterm{def:b1:elementary-steps:catalyst}{catalyst} works only hot, and the conditions are a compromise discussed in
the Year~2 volume.
\textbf{9.} \qty{58.8}{kmol} of ammonia: \qty{29.41}{kmol} of \ce{N2},
\qty{824}{kg}; \qty{88.2}{kmol} of \ce{H2}, \qty{176}{kg}.
\textbf{10.} $V(\ce{N2}) = 29.4 \times 10^3 \times 8.314 \times 298.15/10^5 =
\qty{729}{m^3}$; air $729/0.78 = \qty{935}{m^3}$.
\textbf{11.} From natural gas (steam reforming of methane), with a
\ce{CO2} by-product.
\textbf{12.} Injected into soil it is a cheap fertiliser, but it is a toxic
gas; solids (urea, ammonium nitrate) are easier to store and spread.
\textbf{13.} N: $-$III $\to$ $+$II (5 electrons), O: 0 $\to$ $-$II (4 per
\ce{O2}): \ce{4NH3 + 5O2 -> 4NO + 6H2O}.
\textbf{14.} \ce{2NO + O2 -> 2NO2}.
\textbf{15.} \ce{3NO2 + H2O -> 2HNO3 + NO}: a \omterm{def:b1:e-ph-diagrams:disproportionation}{disproportionation} (N $+$IV to
$+$V and $+$II).
\textbf{16.} \ce{NH3 + 2O2 -> HNO3 + H2O}.
\textbf{17.} $2 \times 58.8 \times 32.0 = \qty{3.76}{t}$.
\textbf{18.} It makes the oxidation fast and selective for NO, instead of
the more stable \ce{N2}.
\textbf{19.} \ce{NH3 + HNO3 -> NH4NO3}.
\textbf{20.} Half.
\textbf{21.} \qty{29.4}{kmol} of \ce{NH4NO3}, \qty{2.35}{t}.
\textbf{22.} $28.0/80.0 = \qty{35}{\percent}$.
\textbf{23.} It contains an \omterm{def:b1:nernst:couple}{oxidant} (nitrate) and a \omterm{def:b1:nernst:couple}{reductant} (ammonium) in
the same \omterm{def:b1:crystals-metals:crystal}{crystal}: heated or mixed with fuels, it can decompose explosively.
\textbf{24.} $58.8 \times 63.0 =$ \textbf{\qty{3.7}{t} of nitric acid per tonne
of ammonia}.
\end{solution}
